\documentclass[12pt,a4paper]{ctexart}

\usepackage[a4paper,margin=2.2cm]{geometry}
\usepackage{graphicx}
\usepackage{float}
\usepackage{booktabs}
\usepackage{array}
\usepackage{hyperref}
\usepackage{xcolor}
\usepackage{enumitem}
\usepackage{caption}
\usepackage{subcaption}
\usepackage{fancyhdr}
\usepackage{lastpage}
\usepackage{listings}
\usepackage{tikz}
\usetikzlibrary{arrows.meta,positioning,shapes.geometric}

\hypersetup{
  colorlinks=true,
  linkcolor=blue,
  urlcolor=blue,
  citecolor=blue
}

\pagestyle{fancy}
\fancyhf{}
\setlength{\headheight}{15pt}
\fancyfoot[C]{第 \thepage 页，共 \pageref{LastPage} 页}

\lstset{
  basicstyle=\ttfamily\small,
  frame=single,
  breaklines=true,
  numbers=left,
  numberstyle=\tiny,
  keywordstyle=\color{blue},
  commentstyle=\color{teal},
  showstringspaces=false
}

% 以下信息建议学生按自己的实际情况直接修改。
\newcommand{\course}{嵌入式系统设计}
\newcommand{\homeworkname}{作业2 LED闪烁实验}
\newcommand{\studentname}{张三}
\newcommand{\studentid}{2024123456}
\newcommand{\classname}{微电子 X 班}
\newcommand{\submitdate}{\today}

\fancyhead[L]{\course}
\fancyhead[R]{\homeworkname}

\newcommand{\placeholderfigure}[2]{
  \begin{figure}[H]
    \centering
    \IfFileExists{#1}{
      \includegraphics[width=0.95\linewidth]{#1}
    }{
      \fbox{\parbox[c][5cm][c]{0.9\linewidth}{\centering
      请将图片文件放入 \texttt{#1}\\
      然后重新编译文档}}
    }
    \caption{#2}
  \end{figure}
}

\begin{document}

\pagenumbering{arabic}
\setcounter{tocdepth}{1}

\begin{center}
  {\zihao{3}\bfseries \homeworkname \par}
  \vspace{0.8em}
  \begin{tabular}{ccc}
    \studentname & \studentid & \classname
  \end{tabular}
  \\
  \vspace{0.4em}
  \submitdate
\end{center}

\vspace{0.5em}
\tableofcontents
\vspace{0.8em}

\section{作业目标与任务要求}

本节简要说明本次作业要完成的功能、验收标准和个人任务分工。建议从以下角度撰写：

\begin{enumerate}[label=\arabic*.]
  \item 本次实验需要实现什么功能。
  \item 使用了哪些硬件资源和软件工具。
  \item 现场验收或结果判定标准是什么。
\end{enumerate}

示例：本实验基于 STM32F103C8T6 完成 GPIO 或 USB 功能验证，要求能够在开发板上稳定运行，并在报告中说明配置过程、核心代码、实验现象及问题分析。

\section{实验环境}

\subsection{硬件环境}

\begin{itemize}
  \item 开发板：STM32F103C8T6 开发板
  \item 调试器：ST-Link
  \item 连接器件：USB 数据线、杜邦线、按键、LED 等
\end{itemize}

\subsection{软件环境}

\begin{itemize}
  \item STM32CubeMX
  \item VS Code 或其他课程指定开发环境
  \item ARM GCC / CMake / Ninja
  \item 串口调试助手、逻辑分析仪软件等
\end{itemize}

\section{实验原理与方案说明}

本节应结合具体作业内容，说明所用外设或中间件的原理。例如：GPIO 输入输出、定时延时、外部中断、USB CDC 通信等。建议从以下三个方面展开：

\begin{enumerate}[label=\arabic*.]
  \item 硬件连接关系。
  \item 外设配置原理。
  \item 程序执行逻辑。
\end{enumerate}

\section{STM32CubeMX 配置过程}

请按实际作业写明关键配置步骤，并至少说明以下内容：

\begin{enumerate}[label=\arabic*.]
  \item 芯片型号选择。
  \item RCC 时钟配置。
  \item SYS 调试接口配置。
  \item GPIO / EXTI / USB 等外设配置。
  \item Clock Configuration 页面中的关键时钟设置。
  \item 工程生成方式与工具链选择。
\end{enumerate}

\placeholderfigure{figures/cubemx-config.png}{CubeMX 关键配置截图}

\section{程序设计与关键代码说明}

本节建议分为“主程序结构”和“关键函数说明”两部分。代码不要大段无解释地堆叠，建议只保留核心片段。

\subsection{主程序结构}

\begin{lstlisting}[language=C,caption={主循环关键代码示例}]
while (1)
{
  // Fill in your main loop logic here.
}
\end{lstlisting}

请在正文中说明：

\begin{itemize}
  \item 每段关键代码实现了什么功能。
  \item 为什么放在 USER CODE 区域。
  \item 若涉及中断、共享变量或通信缓冲区，如何保证数据正确性。
\end{itemize}

\placeholderfigure{figures/code-main.png}{关键代码截图}

\section{程序流程图}

下面给出一个可直接修改的流程图模板。将节点文字替换为本次实验的实际流程即可。

\begin{figure}[H]
  \centering
  \begin{tikzpicture}[
    node distance=10mm and 18mm,
    every node/.style={font=\small},
    startstop/.style={rectangle, rounded corners, minimum width=28mm, minimum height=9mm, draw=black, fill=gray!15},
    process/.style={rectangle, minimum width=34mm, minimum height=9mm, draw=black, fill=blue!8},
    decision/.style={diamond, aspect=2, draw=black, fill=orange!12, inner sep=1pt},
    line/.style={-Latex, thick}
  ]
    \node[startstop] (start) {开始};
    \node[process, below=of start] (init) {系统初始化};
    \node[process, below=of init] (task) {执行主功能};
    \node[decision, below=of task, yshift=-2mm] (judge) {是否满足条件};
    \node[process, below left=of judge] (branch1) {执行分支 A};
    \node[process, below right=of judge] (branch2) {执行分支 B};
    \node[startstop, below=18mm of judge] (endloop) {循环或结束};

    \draw[line] (start) -- (init);
    \draw[line] (init) -- (task);
    \draw[line] (task) -- (judge);
    \draw[line] (judge) -- node[left] {是} (branch1);
    \draw[line] (judge) -- node[right] {否} (branch2);
    \draw[line] (branch1) |- (endloop);
    \draw[line] (branch2) |- (endloop);
  \end{tikzpicture}
  \caption{程序流程图示例}
\end{figure}

\section{实验结果与现象分析}

本节必须体现“实验是否真实完成”。建议至少包含以下内容：

\begin{enumerate}[label=\arabic*.]
  \item 实验现象描述。
  \item 与预期结果是否一致。
  \item 若不一致，可能原因是什么。
\end{enumerate}

\placeholderfigure{figures/result-photo.jpg}{实验现象照片}

\placeholderfigure{figures/result-serial.png}{运行结果截图或上位机交互截图}

\section{问题分析与调试记录}

建议按“问题现象 $\rightarrow$ 排查过程 $\rightarrow$ 原因定位 $\rightarrow$ 解决方法”的格式撰写。例如：

\begin{table}[H]
\centering
\begin{tabular}{p{0.18\textwidth}p{0.74\textwidth}}
\toprule
问题现象 & 解决过程记录 \\
\midrule
程序无法下载 & 检查 SWD 接口、开发板供电、BOOT 配置、调试器驱动等。 \\
实验现象不正确 & 检查引脚配置、时钟配置、代码逻辑和硬件极性。 \\
运行结果不稳定 & 分析抖动、通信忙状态、时序问题或资源冲突。 \\
\bottomrule
\end{tabular}
\caption{问题分析与调试记录示例}
\end{table}

\section{思考题或课后作业回答}

请将本次作业 README 中要求回答的问题逐条作答。建议每题单独编号，回答时尽量结合实验过程，不要只写结论。

\section{总结}

总结本次实验中掌握的关键知识点、存在的问题以及后续可改进的方向。

\section{附录}

可在附录中补充：

\begin{itemize}
  \item 完整关键代码。
  \item 更多截图或照片。
  \item 参考资料。
  \item 自己完成的扩展任务说明。
\end{itemize}

\end{document}